\chapter{FUNDAMENTAÇÃO TEÓRICA}

Este capítulo aborda os conceitos fundamentais que baseiam a arquitetura do CreditMaster, focando no uso de Modelos de Linguagem e técnicas avançadas de engenharia de \textit{prompt} e adaptação de modelos.

\section{Modelos de Linguagem de Grande Escala (LLMs)}

A área de Processamento de Linguagem Natural (NLP) sofreu uma transformação radical com o advento de arquiteturas baseadas em redes neurais profundas (\textit{Deep Learning}). Conforme delineado por \citeonline{lecun2015deeplearning}, o aprendizado profundo permite que modelos computacionais compostos por múltiplas camadas de processamento aprendam representações de dados com múltiplos níveis de abstração. Na seara do NLP, essa evolução culminou na arquitetura \textit{Transformer}, introduzida por \citeonline{vaswani2017attention}, que dispensou os componentes recorrentes e convolucionais baseando-se inteiramente em mecanismos de "atenção". Esse marco permitiu paralelização massiva e o treinamento de Modelos de Linguagem de Grande Escala (LLMs) com bilhões, e eventualmente trilhões, de parâmetros.

Os LLMs são sistemas treinados em vastos \textit{corpora} de texto da internet com o objetivo de prever o próximo \textit{token} (palavra ou sub-palavra) em uma sequência, o que lhes confere notável habilidade de gerar textos coerentes e contextuais. No escopo deste projeto, o modelo \textbf{Llama 3.1 8B Instruct}, desenvolvido de forma \textit{open-source} pela Meta, destaca-se por oferecer uma relação custo-benefício excepcional. Com 8 bilhões de parâmetros, ele atinge um equilíbrio ideal, sendo poderoso o suficiente para assimilar o intrincado jargão financeiro, mas compacto o bastante para ser executado integralmente \textit{on-premise} (localmente) em uma única GPU, evitando gargalos e riscos inerentes ao envio de dados sigilosos para a nuvem.

\section{\textit{Fine-Tuning} e QLoRA}
O ajuste fino (\textit{Fine-Tuning}) completo de modelos com bilhões de parâmetros exige altíssimo poder computacional. Para contornar esse desafio técnico, utiliza-se o método QLoRA (\textit{Quantized Low-Rank Adaptation}), que permite modificar estruturalmente os pesos neurais do modelo base com uso otimizado e drástico de memória. No projeto, a biblioteca \textbf{Unsloth} é utilizada para acelerar esse processo e viabilizar o treinamento em GPUs convencionais.

\section{\textit{Retrieval-Augmented Generation} (RAG)}
A Geração Aumentada por Recuperação (RAG) consiste na injeção dinâmica de contexto factual no momento da inferência do modelo. Para o domínio bancário, a eliminação de "alucinações" (geração de textos plausíveis, porém incorretos ou inventados) é um requisito crítico. A arquitetura RAG permite consultas contra manuais internos e políticas de crédito, garantindo que o modelo responda estritamente de acordo com as regras estabelecidas, mantendo a rastreabilidade do processo decisório.
\textbf{[ANOTAÇÃO: Tenho que expandir essa explicação sobre os bancos vetoriais, falando mais do ChromaDB e FAISS que usei pra salvar e buscar os documentos.]}
